We now prove Theorem~\ref{thm:global-boundary-regularity}.

\begingroup
\renewcommand{\endProof}{\par\nobreak\hfill$\blacksquare$\par\vspace{0.1cm}}
\begin{Proof}
It is now a simple matter to put together the proof of the global regularity theorem. We simply provide an open cover for $\overline\Omega$ using the neighborhoods $\widetilde Q$ constructed in Lemma~\ref{lem:local-boundary-regularity-transfer} for each point $\overline{\bm{x}}\in\partial\Omega$ and one additional set $\Omega_0\subset\subset\Omega$ to cover the interior. Since $\overline\Omega$ is compact, there is a finite subcover (in which we assume $\Omega_0$ is included and which we label $\{\Omega_i\}_{i=0}^N$) such that
\begin{equation}\label{eq:global-boundary-regularity-finite-open-cover}
\Omega\subset\bigcup_{i=0}^N\Omega_i.
\end{equation}
Now, using the interior regularity result (Theorem~\ref{thm:second-order-interior-regularity}) for $\Omega_0$ and Lemma~\ref{lem:local-boundary-regularity-transfer} for each of the other sets we get
\begin{equation}\label{eq:global-boundary-regularity-summed-local-estimates}
\|u\|_{H^2(\Omega)}\leq\sum_{i=0}^N\|u\|_{H^2(\Omega_i)}\leq C(\|f\|_{L^2(\Omega)}+\|u\|_{H^1(\Omega)}).
\end{equation}
A standard application of Ehrling\textquotesingle s lemma (Theorem~\ref{thm:ehrling-compact-intermediate-norm}) gives us the final result.
\end{Proof}
\endgroup
